\section{Timeline Semantics}\label{sec:temporal-model}

Motivo schedules events with a per-context cursor measured in beats.
\topic{Relative placement.}
A \texttt{play} without \texttt{from} places material at the current cursor.
Notes and rests advance the cursor by their duration; chords advance by the longest member; sequences advance
by the sum of their parts.
The colon operator sets an explicit duration (\texttt{play C4:0.5}).
\topic{Absolute placement.}
\texttt{from} schedules at a fixed beat without moving the cursor (\texttt{play A4 from 16;}), so layers can
be added without disturbing sequential flow.
\topic{Control flow and patterns.}
\texttt{loop}, \texttt{for}, and \texttt{if} reuse the enclosing cursor: each loop iteration continues from
the beat reached by the previous one.
Pattern calls run against an internal cursor starting at zero; when the call completes, the caller jumps
forward by the pattern footprint (including internal rests).
\topic{Voices.}
Each voice inherits the parent track cursor or a \texttt{from} offset, then maintains a parallel cursor.
Events from voices merge into the same track output sorted by absolute beat; outer-track statements scheduled
after a voice do not wait for the voice to finish.

\begin{figure}[htbp]
  \centering
  \motivoscalefig{%
    \begin{tikzpicture}[
        beat/.style={motivo/beat},
        note/.style={rounded corners=2pt, minimum height=0.4cm, draw=white, line width=0.8pt},
        arrow/.style={-Stealth, thick, draw=motivoInk!70}
      ]
      \draw[thick, draw=motivoInk!70] (0,0) -- (11.5,0);
      \foreach \x/\b in {0/0,2/1,4/2,6/3,8/4,10/5,11.5/6} {
        \draw[draw=motivoInk!45] (\x,0.12) -- (\x,-0.12);
        \node[beat, below=0.16cm] at (\x,0) {\b};
      }

      \node[note, fill=motivoBlue!70, minimum width=1.7cm] at (1,0.6) {\textcolor{white}{\scriptsize play C4}};
      \node[note, fill=motivoGreen!70, minimum width=1.7cm] at (3,0.6) {\textcolor{white}{\scriptsize play D4}};
      \node[note, fill=motivoAmber!85, minimum width=1.7cm] at (5,0.6) {\textcolor{white}{\scriptsize play E4}};
      \node[note, fill=motivoRose!80, minimum width=2.0cm] at (8,1.35) {\textcolor{white}{\scriptsize play G4 from 4}};

      \draw[arrow] (0,1.05) -- node[beat, above, align=center] {sequential \texttt{play}\\advances cursor} (6.2,1.05);
      \draw[dashed, thick, draw=motivoInk!55] (8,0.9) -- (8,0.5);
      \node[beat, align=center] at (11.2,1.55) {\texttt{from 4}: beat~4 placement,\\cursor unchanged};
    \end{tikzpicture}%
  }
  \caption{Cursor-based temporal placement.
    Sequential \texttt{play} statements advance the cursor; \texttt{from} places an event at an explicit beat without
  moving the cursor.}\label{fig:cursor-model}
\end{figure}
